% Options for packages loaded elsewhere
\PassOptionsToPackage{unicode}{hyperref}
\PassOptionsToPackage{hyphens}{url}
\PassOptionsToPackage{dvipsnames,svgnames,x11names}{xcolor}
%
\documentclass[
  11pt,
]{article}
\usepackage{amsmath,amssymb}
\usepackage{iftex}
\ifPDFTeX
  \usepackage[T1]{fontenc}
  \usepackage[utf8]{inputenc}
  \usepackage{textcomp} % provide euro and other symbols
\else % if luatex or xetex
  \usepackage{unicode-math} % this also loads fontspec
  \defaultfontfeatures{Scale=MatchLowercase}
  \defaultfontfeatures[\rmfamily]{Ligatures=TeX,Scale=1}
\fi
\usepackage{lmodern}
\ifPDFTeX\else
  % xetex/luatex font selection
\fi
% Use upquote if available, for straight quotes in verbatim environments
\IfFileExists{upquote.sty}{\usepackage{upquote}}{}
\IfFileExists{microtype.sty}{% use microtype if available
  \usepackage[]{microtype}
  \UseMicrotypeSet[protrusion]{basicmath} % disable protrusion for tt fonts
}{}
\makeatletter
\@ifundefined{KOMAClassName}{% if non-KOMA class
  \IfFileExists{parskip.sty}{%
    \usepackage{parskip}
  }{% else
    \setlength{\parindent}{0pt}
    \setlength{\parskip}{6pt plus 2pt minus 1pt}}
}{% if KOMA class
  \KOMAoptions{parskip=half}}
\makeatother
\usepackage{xcolor}
\usepackage[margin=1.05in]{geometry}
\setlength{\emergencystretch}{3em} % prevent overfull lines
\providecommand{\tightlist}{%
  \setlength{\itemsep}{0pt}\setlength{\parskip}{0pt}}
\setcounter{secnumdepth}{-\maxdimen} % remove section numbering
\usepackage{amsmath,amssymb,mathtools}
\usepackage{microtype}
\ifLuaTeX
  \usepackage{selnolig}  % disable illegal ligatures
\fi
\usepackage{bookmark}
\IfFileExists{xurl.sty}{\usepackage{xurl}}{} % add URL line breaks if available
\urlstyle{same}
\hypersetup{
  pdftitle={Poisson Approximation and a Sharp Berry--Esseen Theorem for Total Luce Cycles},
  pdfauthor={Christopher D. Long},
  colorlinks=true,
  linkcolor={black},
  filecolor={Maroon},
  citecolor={Blue},
  urlcolor={blue},
  pdfcreator={LaTeX via pandoc}}

\title{Poisson Approximation and a Sharp Berry--Esseen Theorem for Total
Luce Cycles}
\author{Christopher D. Long}
\date{Research note -- 11 October 2026}

\begin{document}
\maketitle

\section{Main result}\label{main-result}

Let \(X_{n,i}\) be independent exponential clocks with positive rates
\(\theta_{n,i}\), and let \(\sigma_n(i)\) be the rank of \(X_{n,i}\).
Write \(C_{n,r}\) for the number of \(r\)-cycles of \(\sigma_n\) and

\[
K_n=\sum_{r=1}^n C_{n,r},\qquad \mu_n=\mathbb E K_n,\qquad L=\log n.
\]

Assume that, after a common rescaling of each row,

\[
1\le\theta_{n,i}\le\kappa
\]

for one fixed \(\kappa<\infty\). No limiting rate profile, smoothness,
or ordering of the rates is assumed.

\textbf{Theorem 1 (Total-cycle Poisson approximation and sharp normal
approximation).} Uniformly over the above rate arrays, for all
sufficiently large \(n\),

\[
\boxed{
 d_{\mathrm{TV}}\!\left(\mathcal L(K_n),\operatorname{Pois}(\mu_n)\right)
 \le C_\kappa\frac{(\log\log n)^2}{\log n}.}
\tag{1}
\]

Moreover,

\[
\mu_n=\log n+O_\kappa(\log\log n),
\tag{2}
\]

and

\[
\boxed{
 d_{\mathrm K}\!\left(
 \mathcal L\!\left(\frac{K_n-\mu_n}{\sqrt{\log n}}\right),
 N(0,1)\right)
 \le\frac{C_\kappa}{\sqrt{\log n}}.}
\tag{3}
\]

The same bound holds with denominator \(\sqrt{\mu_n}\). Here
\(d_{\mathrm{TV}}\) is the supremum over events, equivalently one half
of the \(\ell^1\) distance of probability masses, and \(d_{\mathrm K}\)
is the supremum of the difference of distribution functions.

\textbf{Corollary 2 (Order optimality).} Uniformly over the same arrays,

\[
\max_{k\in\mathbb Z}\mathbb P(K_n=k)
=\frac{1+o_\kappa(1)}{\sqrt{2\pi\log n}}.
\tag{4}
\]

In particular,

\[
 d_{\mathrm K}\!\left(
 \mathcal L\!\left(\frac{K_n-\mu_n}{\sqrt{\log n}}\right),N(0,1)\right)
 \ge
 \left(\frac1{2\sqrt{2\pi}}-o_\kappa(1)\right)(\log n)^{-1/2}.
\tag{5}
\]

Thus the exponent in (3) is optimal. The constant in (5) is a lattice
lower bound, not a claim about the exact leading Kolmogorov constant.

\section{Source and dependency
boundary}\label{source-and-dependency-boundary}

This note develops a new consequence of the quantitative switching
machinery already present in the repository
\texttt{long-mathematics/nonuniform-random-permutations}, frozen at
commit

\texttt{6178a99a8f2ee7b6416edd026e7780b9b8fa9fa0}.

The inputs used are the free-depth endpoint estimate in Paper II, its
polynomial switching specialization and passive-mark Campbell--Stein
argument in Paper IV, and the one- and two-cycle bounds in Paper III.
The precise source labels appear in the final section and the
accompanying source manifest.

Only these modules of Paper IV are needed. Its sharp normal
approximation for logarithmic permutation order, prime-power
linearization, high-frequency arithmetic analysis, and arithmetic
remainder estimates are not inputs to this proof. In particular, the
count theorem does not depend on accepting Paper IV's entire
order-Edgeworth argument.

The new transfer argument is proved below. The upstream endpoint
construction is imported from the previously audited switching theory
rather than re-proved here. Only its annealed one- and two-root
observable estimates are used; no fully quenched arbitrary-deletion
theorem or growing-number-of-roots theorem is assumed.

\section{Quantitative switching
inputs}\label{quantitative-switching-inputs}

For distinct labels let

\[
R_{ij}=\exp\{-(\theta_{n,i}-\theta_{n,j})(X_{n,j}-X_{n,i})\},
\]

and put

\[
S_{n,r}=\sum_{i:L_n(i)>r}R_{i,\sigma_n^r(i)},
\qquad
\Delta_{n,r}=\mathbb E\left|\frac{S_{n,r}}n-1\right|.
\]

Here \(L_n(i)\) is the length of the cycle containing \(i\).

The exact marked switching identity is

\[
 r(n-r)\mathbb E[C_{n,r}\Phi(\sigma_n)]
 =\mathbb E\sum_{i:L_n(i)>r}R_{i,\sigma_n^r(i)}
 \Phi\!\left(\sigma_n\circ(i\ \sigma_n^r(i))\right)
\tag{6}
\]

for bounded permutation-measurable \(\Phi\). The selected-factor bound
gives, for some fixed \(p>1\) depending only on \(\kappa\),

\[
 \sup_{n,r}\frac1n\mathbb E\sum_{i:L_n(i)>r}
 R_{i,\sigma_n^r(i)}^p\le C_{\kappa,p}.
\tag{7}
\]

For \(\kappa>1\), any \(1<p<\kappa/(\kappa-1)\) is admissible; for
\(\kappa=1\), the factors are identically one. Write \(\beta=1-1/p>0\).

The one-cycle and root bounds are

\[
 \mathbb EC_{n,r}\le\frac{C_\kappa n}{r(n-r)},\qquad r<n,
\tag{8}
\]

\[
 \mathbb P\{L_n(I_n)\le a\}\le C_\kappa a/n,
 \qquad a\le n/2,
\tag{9}
\]

where \(I_n\) is an independent uniform root. The two-cycle bound is

\[
 \mathbb E[C_{n,r}(C_{n,s}-\mathbf1_{\{r=s\}})]
 \le\frac{C_\kappa}{rs},\qquad r+s\le n/2.
\tag{10}
\]

Fix a numerical accuracy exponent \(Q\ge4\). Choose \(A=A(\kappa,Q)\)
sufficiently large, and set

\[
 b=\lceil L^3\rceil,\qquad U=\lfloor n/L^A\rfloor,
 \qquad \mathcal I_n=\{r:b<r\le U\},
\tag{11}
\]

\[
 a_{n,r}=\frac{n}{r(n-r)},\qquad
 \lambda_n=\sum_{r\in\mathcal I_n}a_{n,r}.
\tag{12}
\]

The polynomial switching estimate of Paper IV gives

\[
 \sum_j\sup_{2^jb\le r<2^{j+1}b}\Delta_{n,r}
 \le C_{\kappa,Q}L^{-Q},
\tag{13}
\]

where the sum is over blocks with lower endpoint at most \(U\). We
choose \(A\) large enough also that \(A\beta\ge Q+1\).

For completeness, the source of the polynomial accuracy is the free
terminal depth. On a dyadic block with \(\delta\asymp r/n\), take
terminal depth \(m\asymp\log(1/\delta)\). The reset error becomes a
positive power of \(\delta\), while the other terms are logarithmic
factors times positive powers of \(\delta\), or finite-size terms
negligible compared with every fixed power of \(L^{-1}\). The
truncation-uniform one- and two-root observable bounds give

\[
 \Delta_{n,r}\le C_\kappa(\mathcal E_{n,2}+n^{-1})^{1/(2\kappa)}.
\]

Increasing \(A\) therefore makes the one-block error \(O(L^{-Q-2})\);
summing \(O(L)\) blocks proves (13). The dependence on a truncation
level \(T\) is linear for one truncated factor and quadratic for two,
not an uncontrolled fixed-\(T\) constant. No untruncated selected second
moment is required.

\section{Keep the omitted cycles as a passive
mark}\label{keep-the-omitted-cycles-as-a-passive-mark}

Define

\[
 J_n=\sum_{b<r\le U}C_{n,r},\qquad
 D_n=\sum_{r\le b}C_{n,r}+\sum_{r>U}C_{n,r}.
\tag{14}
\]

Thus \(K_n=J_n+D_n\). The variable \(D_n\) is the omitted \textbf{cycle
count}, not the arithmetic remainder for logarithmic order used in Paper
IV.

A good core split has prescribed length \(r\in\mathcal I_n\) and
source-cycle length \(\ell>r+U\). It replaces that source cycle by
cycles of lengths \(r\) and \(\ell-r\). Both the original source and the
residual cycle are longer than \(U\). Hence the split adds one core
cycle, changes no short cycle, and leaves the number of long cycles
unchanged. Therefore

\[
 \xi_n\longmapsto\xi_n+e_r,\qquad D_n\longmapsto D_n,
 \qquad \xi_n=(C_{n,r})_{r\in\mathcal I_n}.
\tag{15}
\]

This is the point at which carrying the full omitted count is stronger
than simply discarding it. There is no assumption that short and long
cycles are independent of one another.

\subsection{The marked approximation}\label{the-marked-approximation}

For arbitrary measurable
\(g_r:\mathbb Z_+^{\mathcal I_n}\times\mathbb Z_+\to[-1,1]\), apply (6)
with

\[
 \Phi_r(\pi)=
 \begin{cases}
 g_r(\xi(\pi)-e_r,D(\pi)),&C_r(\pi)>0,\\
 0,&C_r(\pi)=0.
 \end{cases}
\]

On a good split the test on the right becomes \(g_r(\xi_n,D_n)\). On a
bad split its difference from that value is at most two. Thus

\[
\begin{split}
 &\left|\mathbb E[C_{n,r}g_r(\xi_n-e_r,D_n)]
   -a_{n,r}\mathbb Eg_r(\xi_n,D_n)\right|\\
 &\qquad\le a_{n,r}\Delta_{n,r}
 +\frac{2a_{n,r}}n\mathbb E\sum_{i:r<L_n(i)\le r+U}
 R_{i,\sigma_n^r(i)}.
\end{split}
\tag{16}
\]

By Holder's inequality, (7), and (9), the normalized weighted
exceptional mass is at most \(C_\kappa(U/n)^\beta\). Since
\(a_{n,r}\le2/r\) and each dyadic block has bounded harmonic mass,
summing (16) gives

\[
\begin{split}
 &\left|\mathbb E\sum_{r\in\mathcal I_n}C_{n,r}g_r(\xi_n-e_r,D_n)
   -\sum_{r\in\mathcal I_n}a_{n,r}\mathbb Eg_r(\xi_n,D_n)\right|\\
 &\qquad\le C_\kappa L^{-Q}+C_\kappa L(U/n)^\beta
 \le C_\kappa L^{-Q}.
\end{split}
\tag{17}
\]

The passive-mark Campbell--Stein lemma now gives

\[
 d_{\mathrm{TV}}\!\left(
 \mathcal L(\xi_n,D_n),
 \left(\bigotimes_{r\in\mathcal I_n}\operatorname{Pois}(a_{n,r})\right)
 \otimes\mathcal L(D_n)\right)
 \le C_\kappa L^{-Q}.
\tag{18}
\]

The reference core and the copied mark are independent. To recall why
the Stein constant is dimension-free, solve the immigration--death Stein
equation for an event test with the mark held fixed. Chains whose
initial conditions differ by one core particle can be coupled so that
their difference survives until time \(t\) with probability \(e^{-t}\).
The first difference of the Stein solution is consequently bounded by
\(\int_0^\infty e^{-t}\,dt=1\). Applying (17) proves (18). No new
endpoint theorem for an increasing number of marks is involved.

Let \(Z_n\sim\operatorname{Pois}(\lambda_n)\) and let \(\widetilde D_n\)
have the exact law of \(D_n\), independently of \(Z_n\). Contraction
under summation yields

\[
 d_{\mathrm{TV}}\!\left(
 \mathcal L(K_n),\mathcal L(Z_n+\widetilde D_n)\right)
 \le C_\kappa L^{-Q}.
\tag{19}
\]

\section{A second-moment bound for the omitted
count}\label{a-second-moment-bound-for-the-omitted-count}

Take \(n\) sufficiently large that \(b<U<n/4\). Put

\[
 \mathcal B_n=\{r:r\le b\}\cup\{r:U<r\le n/4\},
 \qquad V_n=\sum_{r\in\mathcal B_n}C_{n,r},
\]

\[
 H_n=\sum_{r\in\mathcal B_n}\frac1r=O_\kappa(\log L).
\]

Every two lengths in \(\mathcal B_n\) sum to at most \(n/2\). Equations
(8) and (10) give

\[
 \mathbb EV_n\le C_\kappa H_n,\qquad
 \mathbb E[V_n(V_n-1)]\le C_\kappa H_n^2.
\tag{20}
\]

There can be at most three cycles longer than \(n/4\): four such cycles
would already contain more than \(n\) vertices. Hence \(D_n=V_n+W_n\)
with \(0\le W_n\le3\), and

\[
 \mathbb ED_n^2\le2\mathbb EV_n^2+18
 \le C_\kappa(1+H_n+H_n^2)
 \le C_\kappa(\log L)^2.
\tag{21}
\]

This argument does not use the two-cycle estimate outside its permitted
range. In particular, it does not attempt to sum a bound with a singular
\(n-r-s\) denominator near full size. Nor does it need fourth or higher
factorial moments.

\section{Exact-mean bookkeeping}\label{exact-mean-bookkeeping}

Total-variation approximation alone does not justify convergence of an
unbounded first moment. We therefore compute the core mean directly from
(6):

\[
 \delta_n:=\mathbb EJ_n-\lambda_n
 =\sum_{r\in\mathcal I_n}a_{n,r}
 \left(\frac{\mathbb ES_{n,r}}n-1\right).
\]

Using (13),

\[
 |\delta_n|\le\sum_{r\in\mathcal I_n}a_{n,r}\Delta_{n,r}
 \le C_\kappa L^{-Q}.
\tag{22}
\]

Write \(\nu_n=\mathbb ED_n\). Since

\[
 a_{n,r}=\frac1r+\frac1{n-r},
\]

harmonic summation gives

\[
 \lambda_n=L-(A+3)\log L+O_\kappa(1),\qquad
 \lambda_n\ge L/2
\tag{23}
\]

for all sufficiently large \(n\). Equation (21) implies
\(\nu_n=O_\kappa(\log L)\). Consequently

\[
 \mu_n=\lambda_n+\nu_n+\delta_n
 =L+O_\kappa(\log L).
\tag{24}
\]

This proves (2), with the expectation calculation independent of any
attempted moment transfer through total variation.

\section{An elementary second-order Poisson smoothing
lemma}\label{an-elementary-second-order-poisson-smoothing-lemma}

\textbf{Lemma 3 (Poisson core plus an independent nonnegative integer
mark).} Let \(Z\sim\operatorname{Pois}(\lambda)\) with \(\lambda>0\),
and let \(M\) be an independent nonnegative integer-valued random
variable with finite second moment. Put \(\nu=\mathbb EM\). Then

\[
\boxed{
 d_{\mathrm{TV}}\!\left(
 \mathcal L(Z+M),\operatorname{Pois}(\lambda+\nu)\right)
 \le
 \frac{\mathbb E[M(M-1)]+\nu^2}{2\sqrt2\,\lambda}.}
\tag{25}
\]

The bound may of course be replaced by its minimum with one.

\textbf{Proof.} Regard the Poisson mass \(p_\lambda\) as a sequence on
all of \(\mathbb Z\), zero on negative integers. Let \(T\) be right
shift, \((Tp)(k)=p(k-1)\), and put \(\Delta=T-I\).

For every integer \(m\ge0\), the exact discrete Taylor identity is

\[
 T^mp=p+m\Delta p+
 \sum_{j=0}^{m-2}(m-1-j)T^j\Delta^2p,
\tag{26}
\]

where the sum is empty for \(m=0,1\). Translation preserves \(\ell^1\)
norm, so

\[
 \|T^mp-p-m\Delta p\|_1
 \le\frac{m(m-1)}2\|\Delta^2p\|_1.
\tag{27}
\]

For a Poisson mass,

\[
 \Delta^2p_\lambda(k)
 =p_\lambda(k)\,
 \frac{(k-\lambda)^2-k}{\lambda^2}.
\tag{28}
\]

The elementary Poisson moment identity

\[
 \mathbb E\left[((Z-\lambda)^2-Z)^2\right]=2\lambda^2
\tag{29}
\]

therefore gives, by Cauchy--Schwarz,

\[
 \|\Delta^2p_\lambda\|_1\le\frac{\sqrt2}{\lambda}.
\tag{30}
\]

Let \(N\sim\operatorname{Pois}(\nu)\), independently of \(Z\). Average
(26) once with \(m=M\) and once with \(m=N\). Both first-order terms are
exactly \(\nu\Delta p_\lambda\), so they cancel. Since
\(\mathbb E[N(N-1)]=\nu^2\),

\[
 \|\mathcal L(Z+M)-\mathcal L(Z+N)\|_1
 \le\frac{\mathbb E[M(M-1)]+\nu^2}{2}
 \|\Delta^2p_\lambda\|_1.
\]

Finally, \(Z+N\sim\operatorname{Pois}(\lambda+\nu)\). Divide by two to
convert the \(\ell^1\) norm into total variation, and use (30). This
proves (25). \(\square\)

The first-moment matching is essential. Directly discarding \(M\) costs
a first-order shift. Replacing it by an independent Poisson variable
with the same mean leaves only a second-order smoothing error. No normal
approximation is needed in this step.

A convenient abstract consequence is the following. If \(K=X+M\) and

\[
 d_{\mathrm{TV}}\bigl(\mathcal L(X,M),
 \operatorname{Pois}(\lambda)\otimes\mathcal L(M)\bigr)\le\varepsilon,
\]

then, with \(\mu=\mathbb EK\) and \(d=\mathbb EX-\lambda\),

\[
 d_{\mathrm{TV}}(\mathcal L(K),\operatorname{Pois}(\mu))
 \le\varepsilon+
 \frac{\mathbb E[M(M-1)]+(\mathbb EM)^2}{2\sqrt2\,\lambda}
 +|d|.
\tag{31}
\]

The last term follows from the elementary Poisson parameter coupling
\(d_{\mathrm{TV}}(\operatorname{Pois}(u),\operatorname{Pois}(v))\le |u-v|\).

\section{Proof of the total-cycle Poisson
approximation}\label{proof-of-the-total-cycle-poisson-approximation}

Apply Lemma 3 with \(Z=Z_n\) and \(M=\widetilde D_n\). By (21) and (23),

\[
 d_{\mathrm{TV}}\!\left(
 \mathcal L(Z_n+\widetilde D_n),
 \operatorname{Pois}(\lambda_n+\nu_n)\right)
 \le C_\kappa\frac{(\log L)^2}{L}.
\tag{32}
\]

The difference between \(\lambda_n+\nu_n\) and the exact mean \(\mu_n\)
is \(\delta_n\), bounded separately in (22). Combining (19), (22), and
(32) proves

\[
 d_{\mathrm{TV}}\!\left(\mathcal L(K_n),\operatorname{Pois}(\mu_n)\right)
 \le C_\kappa\left[L^{-Q}+\frac{(\log L)^2}{L}\right].
\tag{33}
\]

Since \(Q\ge4\) is fixed and \(A\) has already been chosen in terms of
\(\kappa\), this is (1). In particular, the Poisson approximation error
is

\[
 O_\kappa((\log L)^2/L)=o_\kappa(L^{-1/2}).
\tag{34}
\]

\section{Proof of the sharp Berry--Esseen
bound}\label{proof-of-the-sharp-berryesseen-bound}

Let \(P_n\sim\operatorname{Pois}(\mu_n)\). The ordinary Poisson normal
approximation gives

\[
 d_{\mathrm K}\!\left(
 \mathcal L\!\left(\frac{P_n-\mu_n}{\sqrt{\mu_n}}\right),N(0,1)\right)
 \le C\mu_n^{-1/2}.
\tag{35}
\]

For example, (35) follows directly from the classical Berry--Esseen
theorem by expressing \(P_n\) as the sum of \(N=\lceil\mu_n\rceil\)
independent Poisson variables of mean \(\mu_n/N\le1\). Their total third
absolute centered moment is \(O(\mu_n)\) and their total variance is
\(\mu_n\). No regularity or nonlattice assumption is involved.

Contraction of total variation under affine maps and (33) prove (3) with
scale \(\sqrt{\mu_n}\). To replace it by \(\sqrt L\), use (24):

\[
 \frac{\mu_n}{L}=1+O_\kappa(\log L/L).
\]

The Kolmogorov distance between centered normals with this variance
ratio is \(O_\kappa(\log L/L)\). Hence

\[
\begin{split}
 d_{\mathrm K}\!\left(
 \mathcal L\!\left(\frac{K_n-\mu_n}{\sqrt L}\right),N(0,1)\right)
 &\le C_\kappa\left[
 L^{-1/2}+\frac{(\log L)^2}{L}+\frac{\log L}{L}\right]\\
 &\le C_\kappa L^{-1/2}.
\end{split}
\tag{36}
\]

This proves Theorem 1. The main normal approximation error comes from
the Poisson law itself, not from restoring the omitted cycles.

\section{Lattice optimality and a local
consequence}\label{lattice-optimality-and-a-local-consequence}

A Poisson distribution of mean \(\mu\to\infty\) has maximal atom

\[
 \max_k\mathbb P\{\operatorname{Pois}(\mu)=k\}
 =\frac{1+o(1)}{\sqrt{2\pi\mu}},
\]

by Stirling's formula. Its mass at \(k=\lfloor\mu\rfloor\) has this same
asymptotic. Equation (33) bounds every individual mass error by
\(o_\kappa(L^{-1/2})\), while \(\mu_n/L\to1\) uniformly. The upper bound
on all atoms and the lower bound at \(\lfloor\mu_n\rfloor\) prove (4).

If a distribution function has a jump of size \(p\) and the comparison
distribution function is continuous, their Kolmogorov distance is at
least \(p/2\). Indeed, at the jump point and its left limit, continuity
of the comparison function prevents both errors from being smaller than
\(p/2\). Applying this observation with (4) proves (5).

The same argument, together with the uniform local Poisson limit, also
yields

\[
 \sup_{k\in\mathbb Z}
 \left|\sqrt L\,\mathbb P(K_n=k)
 -\varphi\!\left(\frac{k-\mu_n}{\sqrt L}\right)\right|
 \longrightarrow0
\tag{37}
\]

uniformly over each fixed comparability class. This local consequence is
not needed for the Berry--Esseen theorem.

\section{What this does not assert}\label{what-this-does-not-assert}

The mean estimate (2) is not the sharper profile-dependent expansion
\(\mathbb EK_n=\log n+\kappa_f+o(1)\). It also does not establish
\(\mu_n-\log n=O_\kappa(1)\) for arbitrary comparable arrays.
Consequently, replacing the exact mean in (3) by \(\log n\) is not
justified at the sharp rate by this proof. The direct bound obtained
with that center is only \(O_\kappa(\log\log n/\sqrt{\log n})\). Any
deterministic center known to be within \(O_\kappa(1)\) of \(\mu_n\)
would preserve the sharp rate.

Neither variance convergence nor higher-cumulant asymptotics are
inferred from total variation. The denominators in the theorem are
\(\sqrt{\log n}\) and \(\sqrt{\mu_n}\), not an unproved identification
with \(\sqrt{\operatorname{Var}K_n}\). An exact-variance formulation can
be treated separately using additional moment control.

A first Edgeworth expansion for \(K_n\) must account for its integer
lattice. A smooth correction copied from the log-order problem cannot
have a uniform \(o(L^{-1/2})\) remainder without an appropriate lattice
correction. No such expansion or exact leading Kolmogorov constant is
claimed in this note.

No repository source has been changed or promoted by this note. The new
proof is available for a separate audit and integration. The general
smoothing lemma is elementary Poisson/Charlier analysis; no claim of
novelty is made for that lemma itself.

\section{Reproducibility and source
map}\label{reproducibility-and-source-map}

The accompanying script \texttt{check\_poisson\_smoothing.py} performs
exact rational checks of the discrete Taylor and Charlier moment
identities, exact finite-permutation checks of the passive-mark split,
and numerical checks of the Poisson convolution bound.

The executed checks passed: 104 discrete Taylor cases; eight rational
Charlier moment cases; 84,438 eligible splits across 5,910 permutations
of sizes three through seven; and 64 numerical convolution cases.
Numerical cases include deterministic, Bernoulli, uniform, binomial, and
rare-large-mark examples. They are consistency checks, not a substitute
for the analytic proof.

\textbf{Repository sources.} The source manifest pins the following
documents at the commit stated above.

\begin{itemize}
\tightlist
\item
  \textbf{{[}II{]}} Christopher D. Long, \emph{Mesoscopic cycles of
  comparable-rate Luce permutations: a switching proof of Conjecture
  13.1}. The exact marked switching identity, selected-factor moments,
  crude root bound, and free-depth one-/two-root endpoint theorem.
\item
  \textbf{{[}III{]}} Christopher D. Long, \emph{An Erdos--Turan Law for
  Comparable-Rate Luce Permutations}. Labels \texttt{input:one-switch},
  \texttt{lem:crude}, and \texttt{cor:pair}. The pair estimate used here
  is explicitly restricted to \(r+s\le n/2\).
\item
  \textbf{{[}IV{]}} Christopher D. Long, \emph{A Sharp Berry--Esseen
  Theorem and First Edgeworth Expansion for the Erdos--Turan Law of
  Comparable-Rate Luce Permutations}. Labels \texttt{input:free-depth},
  \texttt{prop:poly-block}, \texttt{lem:passive-Stein}, and
  \texttt{thm:poly-pp}; in particular source lines 1750--2007. Blob SHA:
  \texttt{28ba5778ecb8bb83da97872e5f68a99401e6b5ce}.
\end{itemize}

\textbf{External background.} V. Yu. Korolev and I. G. Shevtsova,
\emph{An improvement of the Berry--Esseen inequality with applications
to Poisson and mixed Poisson random sums}, arXiv:0912.2795. A. D.
Barbour, E. Kowalski and A. Nikeghbali, \emph{Mod-discrete expansions},
arXiv:0912.1886. These situate the standard normal-approximation and
Poisson/Charlier tools; the total-cycle transfer above is derived from
the stated repository inputs, not from an assumption of mod-Poisson
convergence.

\end{document}
